%% kapa: exam=1 round=35 year=2024 period=2 form=A subject=accounting subject_ko=회계학 seq=17 number=57
%% source: Q-Net 감정평가사 기출문제 (cst00309) artlSeq=5213229 file=raw/1008/5213229/2교시 문제지 원본.pdf
%% conversion: kordoc 4.24.51 -> tools/gen1.py (auto); review=proofread-sonnet; flags=tab
\begin{question}{57}
(주)감평은 20×1년 초 토지 A(취득원가 ￦1,000)와 토지 B(취득원가 ￦2,000)를 각각 취득하고, 재평가모형을 적용하였다. 동 2건의 토지에 대하여 공정가치가 다음과 같을 때, 각 연도별 당기순이익 또는 기타포괄이익에 미치는 영향으로 옳은 것은? (단, 토지에 대한 재평가잉여금의 일부를 이익잉여금으로 대체하지 않는다.)
\par{\small\centering\begin{tabular}{lrrr}\hline
 & \multicolumn{1}{c}{20×1년 말} & \multicolumn{1}{c}{20×2년 말} & \multicolumn{1}{c}{20×3년 말} \tabularnewline\hline
토지 A & ￦1,100 & ￦950 & ￦920 \tabularnewline
토지 B & 1,700 & 2,000 & 2,100 \tabularnewline\hline
\end{tabular}\par}
\begin{choices}
  \item 20×1년 말 토지 A로부터 당기순이익 ￦100이 증가한다.
  \item 20×2년 말 토지 A로부터 당기순이익 ￦150이 감소한다.
  \item 20×2년 말 토지 B로부터 기타포괄이익 ￦300이 증가한다.
  \item 20×3년 말 토지 A로부터 기타포괄이익 ￦30이 감소한다.
  \item 20×3년 말 토지 B로부터 기타포괄이익 ￦100이 증가한다.
\end{choices}
\end{question}
\answer{5}
